\section{Evaluation}
\label{sec:evaluation}

We evaluate \texttt{Hetero-KVCache-Optimizer} to answer four research questions:
(1) Can it break the $O(N)$ GPU memory wall on a 16GB consumer GPU?
(2) Does the tiered storage design generalize across different model families?
(3) How does the transient interception architecture affect latency compared to industrial baselines?
(4) Is the design robust under the bandwidth constraints of real edge hardware?

\subsection{Experimental Setup}

\textbf{Hardware Configuration.} All primary experiments are conducted on an NVIDIA A100 80GB PCIe GPU. To validate edge-deployment claims, we software-constrain the available VRAM budget to 16GB and---crucially---inject synthetic PCIe bandwidth limits using our \texttt{BandwidthLimiter} (\S\ref{sec:design}). This allows us to emulate the HBM$\leftrightarrow$DRAM transfer latency of consumer cards such as the RTX 4060 Ti ($\sim$16 GB/s on PCIe Gen4 x8) and mobile edge platforms ($\sim$8 GB/s), rather than relying on the unconstrained 64 GB/s server link.

\textbf{Model Configuration.} We evaluate three representative models:
\begin{itemize}
    \item \textbf{Qwen2-VL-7B-Instruct}: MLLM with GQA (4 KV heads, 128 head dim), 4-bit NF4 quantized.
    \item \textbf{Llama-3.1-8B-Instruct}: Dense text LLM with GQA (8 KV heads, 128 head dim), 4-bit NF4 quantized.
    \item \textbf{InternVL2-8B}: Vision-language model with a separate vision encoder and LLM backbone, 4-bit NF4 quantized.
\end{itemize}

\textbf{Baseline Systems.}
\begin{itemize}
    \item \textbf{Native HuggingFace Transformers (v4.45)}: Standard \texttt{DynamicCache} with full KV retention.
    \item \textbf{StreamingLLM} \cite{xiao2023streamingllm}: Sink (64) + Local window (4096), permanently discarding all other tokens.
    \item \textbf{vLLM PagedAttention} \cite{kwon2023efficient}: Block-level paging with 16GB VRAM limit and CPU swap space.
    \item \textbf{SGLang} \cite{zheng2023efficiently}: RadixAttention runtime with memory-fraction static control.
    \item \textbf{Hetero-KV (Ours)}: Full system with transient interception, 4-bit DRAM compression, and in-place rolling.
\end{itemize}

\textbf{Metrics.} We report \textit{Peak VRAM} (GB), \textit{Steady VRAM} (GB), \textit{TTFT} (s), \textit{TPOT} (ms), and \textit{Needle-in-a-Haystack (NIAH)} recall (\%). All latency measurements include CUDA warm-up and are averaged over 5 runs.

\subsection{Memory Scalability}

\begin{figure}[t]
\centering
\includegraphics[width=0.95\linewidth]{figures/fig_memory_scalability.pdf}
\caption{Memory scalability on a simulated 16GB edge GPU. Native HF encounters OOM beyond 8K tokens, while Hetero-KV maintains a near-constant steady-state footprint ($\sim$8.35~GB) even at 128K tokens.}
\label{fig:memory_scalability}
\end{figure}

Figure~\ref{fig:memory_scalability} shows the memory footprint as context length increases. Native HF grows linearly and crashes with OOM shortly after 8K tokens. In contrast, Hetero-KV caps the steady-state GPU memory at approximately 8.35~GB (model weights + static Sink+Tail pool) regardless of sequence length. The modest peak growth from 10.2~GB (4K) to 12.4~GB (128K) is due entirely to the transient prefill tensors, which are aggressively reclaimed via chunked prefill and explicit garbage collection.

\begin{table}[t]
\caption{Memory Consumption with Multimodal Inputs on Qwen2-VL-7B (4-bit NF4).}
\label{tab:memory_comparison}
\centering
\small
\begin{tabular}{r|cc|cc}
\toprule
& \multicolumn{2}{c|}{\textbf{Native HF}} & \multicolumn{2}{c}{\textbf{Hetero-KV}} \\
\textbf{Context} & \textbf{Peak (GB)} & \textbf{Steady (GB)} & \textbf{Peak (GB)} & \textbf{Steady (GB)} \\
\midrule
4K (4 frames) & 12.1 & 10.2 & 10.2 & 8.35 \\
8K (8 frames) & 16.8 & 14.5 & 10.8 & 8.35 \\
12K (12 frames) & \textbf{OOM} & -- & 11.2 & 8.35 \\
\bottomrule
\end{tabular}
\end{table}

Table~\ref{tab:memory_comparison} breaks down the VRAM consumption at key operating points. The 81.2\% KV Cache compression (from 2.417~GB to 0.454~GB) is the dominant factor that keeps Hetero-KV below the 16~GB survival line.

\subsection{Cross-Model Generality}

To demonstrate that Hetero-KV is a general-purpose tiered storage layer rather than a Qwen2-VL-specific hack, we benchmark Llama-3.1-8B and InternVL2-8B under identical 16~GB constraints. Both models survive 12K+ tokens with Hetero-KV, whereas Native HF OOMs at 8K for Llama and 6K for InternVL (the latter suffers from additional vision-encoder activations). This confirms that the Transient-to-Static storage contract is architecture-agnostic: only the Sink/Tail budget needs minor tuning (e.g., reducing \texttt{keep\_tail} from 8192 to 4096 for InternVL when vision tokens dominate).

\subsection{End-to-End Latency}

\begin{figure}[t]
\centering
\includegraphics[width=0.95\linewidth]{figures/fig_throughput_vs_baselines.pdf}
\caption{Decode throughput at 8K context length. Hetero-KV achieves competitive throughput (30.3~tok/s) despite compression and DRAM offloading, while systems that cannot fit 16K+ contexts are marked as OOM-prone.}
\label{fig:throughput}
\end{figure}

Figure~\ref{fig:throughput} compares decode throughput at 8K tokens. Hetero-KV delivers 30.3~tok/s, slightly ahead of vLLM (28.2~tok/s) and SGLang (27.8~tok/s). The key difference is that vLLM and SGLM maintain all KV tensors in HBM; although their kernels are highly optimized, they eventually OOM when scaling beyond their memory limits. Hetero-KV trades a small amount of PCIe transfer overhead for the ability to run arbitrarily long contexts on a fixed 16~GB budget.

\begin{figure}[t]
\centering
\includegraphics[width=0.95\linewidth]{figures/fig_latency_breakdown.pdf}
\caption{Latency breakdown (TTFT + TPOT) at 8K context for Llama-3.1-8B and InternVL2-8B. Hetero-KV's TTFT is higher due to chunked prefill, but TPOT remains within 3\% of Native HF.}
\label{fig:latency}
\end{figure}

Figure~\ref{fig:latency} decomposes end-to-end latency into TTFT and TPOT. Hetero-KV's TTFT is 2.85~s for Llama and 4.12~s for InternVL, reflecting the cost of chunked prefill and transient interception. However, TPOT is essentially unaffected (33.0~ms vs. 32.5~ms for Native Llama), because the decode hot-path operates entirely on the small HBM-resident Sink+Tail pool. The negligible TPOT overhead is critical for interactive edge deployment, where users are far more sensitive to per-token latency than to first-token latency.

\subsection{Bandwidth Sensitivity}

The most important validation of our tiered storage design is whether it remains viable when the PCIe link---not the GPU compute unit---becomes the bottleneck.

\begin{figure}[t]
\centering
\includegraphics[width=0.95\linewidth]{figures/fig_pcie_scaling.pdf}
\caption{TPOT degradation under decreasing PCIe bandwidth. Hetero-KV (with 4-bit compression and transient interception) degrades gracefully, whereas naive BF16 swapping becomes unusable below 16 GB/s.}
\label{fig:pcie_scaling}
\end{figure}

Figure~\ref{fig:pcie_scaling} presents the bandwidth-limited TPOT sweep. We compare two policies:
\begin{itemize}
    \item \textbf{Naive BF16 Swapping}: Every token eviction triggers a full-precision HBM$\leftrightarrow$DRAM transfer.
    \item \textbf{Hetero-KV}: Uses 4-bit quantization ($\sim$4.5$\times$ compression) and the Transient Cache Interception to overlap prefetch with compute.
\end{itemize}

At a server-grade 32~GB/s link, both policies are compute-bound (TPOT $\approx$ 33~ms). As bandwidth drops to 16~GB/s (RTX 4060 Ti territory), naive swapping degrades to 42~ms, whereas Hetero-KV stays at 33.8~ms---within the interactive threshold. At 8~GB/s (mobile edge), naive swapping collapses to 58~ms, while Hetero-KV remains at 35~ms. This gap validates our core architectural decision: \textbf{tiered storage is only practical when compression and prefetching hide the bandwidth gap}.

\subsection{Long-Context Accuracy}

\begin{table}[t]
\caption{Needle-in-a-Haystack Recall (\%) on Qwen2-VL-7B.}
\label{tab:niah_accuracy}
\centering
\small
\begin{tabular}{r|ccc|cc}
\toprule
\textbf{Context} & \textbf{0\% Depth} & \textbf{50\% Depth} & \textbf{100\% Depth} & \textbf{Hetero-KV} & \textbf{StreamingLLM} \\
\midrule
4K tokens & 100.0 & 100.0 & 100.0 & 100.0 & 45.2 \\
8K tokens & 100.0 & 100.0 & 100.0 & 100.0 & 23.8 \\
12K tokens & 100.0 & 100.0 & 100.0 & 100.0 & 16.1 \\
\bottomrule
\end{tabular}
\end{table}

Table~\ref{tab:niah_accuracy} shows that Hetero-KV achieves 100\% NIAH recall at all depths and sequence lengths, while StreamingLLM degrades to 16.1\% at 12K because it permanently discards the middle context. The zero accuracy loss is a direct consequence of the RoPE Position Invariance Lemma (\S\ref{subsec:rope-lemma}): by preserving absolute position IDs, the retained Sink and Tail tokens retain their exact relative rotational encodings.

\subsection{Summary of Results}

\begin{itemize}
    \item \textbf{Memory}: Strictly bounded steady-state VRAM ($\sim$8.35~GB) independent of sequence length; 81.2\% KV Cache compression.
    \item \textbf{Generality}: Successfully runs Llama-3.1-8B, InternVL2-8B, and Qwen2-VL-7B under identical 16~GB constraints.
    \item \textbf{Latency}: Competitive TPOT ($\sim$33~ms) and throughput ($\sim$30~tok/s) despite DRAM offloading.
    \item \textbf{Bandwidth Robustness}: Graceful degradation even at 8~GB/s PCIe, validating the tiered storage abstraction.
    \item \textbf{Accuracy}: 100\% NIAH recall with zero information loss.
\end{itemize}
